% HISTORICAL INTEGRATION SNIPPET: promoted on 29 August 2026 after the
% exact-byte production Rethlas pass for the P06.1 statement and proof.

% --------------------------------------------------------------------------
% Chapter 18 hypothesis ledger: insert before the row headed
% ``real-PL zero-current converse''.
% --------------------------------------------------------------------------
quasi-smooth classical pluriharmonic regularity
 &\Cref{thm:p061-quasismooth-ph-rpl}: \(K\) is complete and nontrivially
  valued, and \(X\) is good, strict, and quasi-smooth; goodness and strictness
  provide strict affinoid neighbourhoods and the rigid--Berkovich
  translation, while quasi-smoothness is the smoothness input in
  \cite[Corollary~3.3.2]{Temkin17}; no global quasicompactness,
  separatedness, or boundarylessness is used; non-quasi-smooth points of
  dimension at least two remain outside the conclusion\\

% --------------------------------------------------------------------------
% Chapter 18 comparison ledger: insert immediately before the row
% ``$\mathbf R$-PL function with zero CLD current''.
% --------------------------------------------------------------------------
Classically pluriharmonic function & \(\mathbf R\)-PL and GJR harmonic &
good strict quasi-smooth analytic space over a complete nontrivially valued
field; \Cref{thm:p061-quasismooth-ph-rpl}; no global quasicompactness,
separatedness, or boundarylessness; \(\mathbf Q\)-PL already fails on a
strict semistable annulus by
\Cref{prop:p061-integral-pl-counterexample} \\

% --------------------------------------------------------------------------
% Chapter 18 source notes: insert after the paragraph about the real-PL
% zero-current converse and before the P06.2 proper-moment paragraph.
% --------------------------------------------------------------------------
The quasi-smooth classical-pluriharmonic regularity theorem
\Cref{thm:p061-quasismooth-ph-rpl} combines the strictly semistable case
\cite[Proposition~7.8]{GR25Classical} with altered local uniformization
\cite[Corollary~3.3.2]{Temkin17}.  Its additional bridge is rigid
\(\acute{e}\)tale descent for \(\mathbf R\)-PL functions: the proof applies
the sheaf criterion \cite[Proposition~3.2.2]{dJvdP96}, the finite-flat
descent theorem \cite[Proposition~5.11]{GJR26} only to finite
\(\acute{e}\)tale affinoid morphisms, and ground-field descent
\cite[Corollary~5.12]{GJR26}.  The annular example
\Cref{prop:p061-integral-pl-counterexample} records the sharp coefficient
  firewall.  The exact statement and proof passed production run
  \nolinkurl{20260829T043457Z_0b1b4d161e3b}, item
  \nolinkurl{pi_7cd35ebd843cf80a7cd696bb}, final \texttt{correct}, zero
  critical errors, and zero gaps; no novelty or priority claim is made.
